\documentclass[11pt]{article}
\usepackage[margin=0.8in]{geometry}
\usepackage{booktabs}
\title{Day 89-B: Deterministic Large-Particle Fixture}
\author{VSEPR-SIM Work Order Authority}
\date{Day 89}
\begin{document}
\maketitle
\section{Authority}
This document governs the Day 89-B fixture implemented in \texttt{apps/cg\_anim\_demo.cpp}. The existing four-scene default remains unchanged. The optional \texttt{--stress-particles N} argument appends a deterministic cubic-lattice scene with exactly $N$ rows; \texttt{--steps N} controls environment evolution and defaults to one for stress runs.

\section{Determinism and scale}
The stress lattice is generated by the existing deterministic lattice builder and truncated to the requested count. Particle order, positions, scene identifier, and output row order are stable for identical arguments. The accepted run requested 1,331 stress particles and emitted 1,442 total rows: 111 baseline rows plus 1,331 stress rows.

\begin{center}
\begin{tabular}{lr}
\toprule
Measure & Accepted result \\
\midrule
Stress-scene particles & 1,331 \\
Total artifact particles & 1,442 \\
Scene count & 5 \\
Environment steps & 1 \\
Measured generation & 41.2803 ms \\
\bottomrule
\end{tabular}
\end{center}

\section{Manifest evidence}
The manifest now records \texttt{particle\_count} and \texttt{generation\_ms} in addition to scene count, steps, and time step. This makes scale explicit and prevents report or renderer validation from inferring particle count by sampling a file.

\newpage
\section{Acceptance command}
From the repository root, build the \texttt{dual-backend-3d-data} target with the release preset, then execute \texttt{build/dual-backend-3d-data.exe --headless --output out/day89\_stress --stress-particles 1331}. Parse \texttt{manifest.json}; acceptance requires five scenes, 1,442 total particles, and a positive generation duration.

\section{Complexity boundary}
The environment evaluator constructs per-bead neighborhoods and therefore is not treated as a renderer stress loop. Stress mode defaults to one simulation step so Day 89 can first establish artifact and consumer behavior above 1,000 particles. Longer scientific stress runs remain explicit through \texttt{--steps}; no silent reduction occurs when that option is supplied.

\section{Exit decision}
Day 89-B is complete because the fixture is configurable, deterministic, compiled, and measured above the required scale. Day 89-C owns the expanded renderer-neutral artifact contract. W--Z remain reserved.
\end{document}
